\documentclass[cn,hazy,blue,14pt,screen]{elegantnote}
\title{}
\author{}
\usepackage{array}
\usepackage{amsmath,booktabs}
\usepackage{amssymb}

\begin{document}
\maketitle

\section{消费者均衡公式推导}

$\begin{cases}
    P_x \cdot x + P_y \cdot y = I \,\, (I=Income)\\
    U(x,y) = x^{\alpha} \cdot y^{\beta}
\end{cases} $

$ \frac{MU_x}{MU_y} = \frac{P_x}{P_y} (1)$
$P_x \cdot x + P_y \cdot y = I (2)$

$ MU_x = \frac{\partial U(x,y)}{\partial x} = \alpha x^{\alpha - 1}y^{\beta} $
$ MU_y = \frac{\partial U(x,y)}{\partial y} = \beta x^{\alpha}y^{\beta - 1} $

$\Rightarrow \frac{MU_x}{MU_y} =  \frac{\alpha x^{\alpha-1} y^{\beta}}{\beta x^{\alpha}y^{\beta - 1}} = \frac{\alpha y}{\beta x} $

$\Rightarrow \frac{\alpha y}{\beta x} = \frac{P_x}{P_y} $

$\Rightarrow P_x \cdot x = P_y \cdot y \cdot \frac{\alpha}{\beta} $

$\Rightarrow \frac{\alpha}{\beta}P_y\cdot y + P_y \cdot y = I $

$\Rightarrow (1+\frac{\alpha}{\beta})P_y\cdot y  = I $

$\Rightarrow y= \frac{I}{P_y} \cdot \frac{\beta}{\alpha+\beta}$

同理：

$\Rightarrow (1+\frac{\beta}{\alpha})P_x\cdot x  = I $

$\Rightarrow x= \frac{I}{P_x} \cdot \frac{\alpha}{\alpha+\beta}$


$MRS_{1,2} = \left. -\frac{\Delta Q_2}{\Delta Q_1} \right|_{U\text{不变}}$

\subsection{Example 1}

Given utility function $U = X^{0.5}Y^{0.5}$, income $I = 80$ and $P_x = 2$, $P_y =4$.

Question 1: What is the maximum utility quantity of goods $X$ and $Y$?


Question 2: What is the maximum utility $U$?


Answer(Version 1):
$$U = X^{0.5}Y^{0.5}$$
$$MU_x = \frac{\partial U(x,y)}{\partial x} = 0.5X^{-0.5}Y^{0.5}$$
$$MU_y = \frac{\partial U(x,y)}{\partial y} = 0.5X^{0.5}Y^{-0.5}$$
$$\text{Condition: } \frac{MU_x}{P_x} = \frac{MU_y}{P_y} \Rightarrow \frac{Y}{2} = \frac{X}{4}$$
$$ \text{Budget:} P_x \cdot x + P_y \cdot y = I$$
$$ 2X+4Y=80$$
$$\Rightarrow 2X + 4 \cdot \frac{1}{2} X =80$$
$$ \Rightarrow X=20 \,\, Y=10$$
$$U_{max} = \sqrt{20\times10} = \sqrt{200} = 10\sqrt{2}$$


Answer(Version 2):


$$\begin{cases}
    X= \frac{I}{P_x} \cdot \frac{\alpha}{\alpha+\beta} \\
    Y= \frac{I}{P_y} \cdot \frac{\beta}{\alpha+\beta}   
\end{cases}$$

$$X = \frac{80}{2} \cdot \frac{0.5}{0.5+0.5} = 20$$
$$Y = \frac{80}{4} \cdot \frac{0.5}{0.5+0.5} = 10$$

$$U_{max} = \sqrt{20\times10} = \sqrt{200} = 10\sqrt{2}$$


\section{最优消费组合类型题目}

\subsection{柯布-道格拉斯函数 Cobb‑Douglas}
形如$U = A X^{\alpha}Y^{\beta}, \alpha>0, \beta>0 $ \\
边际效用:$MU_X = \alpha X^{\alpha-1} Y^{\beta}, MU_Y = \beta X^{\alpha}Y^{\beta -1}$ \\
均衡条件:$\frac{MU_X}{MU_Y} = \frac{P_x}{P_y}$ 或者写为 $\frac{MU_X}{P_X} = \frac{MU_Y}{P_Y}$\\

\textbf{Example1:}
消费者的效用函数为$U=X^{0.5}Y^{0.5}$，收入为120元，$P_X=3$，$P_Y=6$。最优消费组合是：

\textbf{解：}

\[
U = X^{0.5}Y^{0.5}
\]

\[
MU_X = \frac{\partial U}{\partial X} = 0.5 X^{-0.5}Y^{0.5},\qquad
MU_Y = \frac{\partial U}{\partial Y} = 0.5 X^{0.5}Y^{-0.5}
\]

最优条件为
\[
\frac{MU_X}{MU_Y} = \frac{P_X}{P_Y}
\]

代入得
\[
\frac{0.5 X^{-0.5}Y^{0.5}}{0.5 X^{0.5}Y^{-0.5}} = \frac{Y}{X} = \frac{3}{6} = \frac12
\]

故
\[
X = 2Y
\]

预算约束为
\[
P_X X + P_Y Y = I
\]
\[
3X + 6Y = 120
\]

把 $X=2Y$ 代入：
\[
3(2Y)+6Y=120 \Rightarrow 12Y=120 \Rightarrow Y=10
\]
\[
X=2Y=20
\]

因此最优消费组合为
\[
X^*=20,\qquad Y^*=10
\]

验证：
\[
3\times 20 + 6\times 10 = 60+60 = 120
\]

故该组合刚好满足预算约束，且为最优消费组合。

\[
U_{\max}=20^{0.5}10^{0.5}=\sqrt{200}=10\sqrt{2}
\]

\subsection{完全替代品 Perfect Subsititutites}
$ U =aX+bY, a>0,b>0$\\
边际效用:$ MU_X = a, MU_Y = b$ \\
边际替代率:$ MRS_{XY} = \frac{MU_X}{MU_Y} = \frac{a}{b}$ \\
无差异曲线：直线\\
均衡规则：\\
1. $\frac{MU_X}{P_X} > \frac{MU_Y}{P_Y}$:全部买$X$，$Y=0$.\\
2. $\frac{MU_X}{P_X} < \frac{MU_Y}{P_Y}$:全部买$Y$，$X=0$.\\
3. $\frac{MU_X}{P_X} = \frac{MU_Y}{P_Y}$:任意组合都是最优的，$P_x \cdot x + P_y \cdot y = I$.\\
\textbf{对于完全替代品不能直接用$\frac{MU_X}{P_X} = \frac{MU_Y}{P_Y}$解方程，否则会出现无限多组解。}

\textbf{Example 1:}某消费者的效用函数为 $U=2X+Y$，收入为120元，$P_X=4$，$P_Y=3$。其最优消费组合是：\\
对于完全替代品，比较的是
\[
\frac{MU_X}{P_X} \quad \text{和} \quad \frac{MU_Y}{P_Y}
\]

\[
MU_X = 2,\qquad MU_Y = 1
\]
\[
\frac{MU_X}{P_X} = \frac{2}{4} = 0.5,\qquad \frac{MU_Y}{P_Y} = \frac{1}{3} \approx 0.3333
\]
\[
\frac{MU_X}{P_X} > \frac{MU_Y}{P_Y}
\]
因此，每一元钱用于购买 $X$ 能带来更高的效用，所以消费者会把全部收入用于购买 $X$，即 $Y=0$。

预算约束为
\[
P_X X + P_Y Y = 120
\]
\[
4X + 3Y = 120
\]
当 $Y=0$ 时，
\[
4X = 120 \Rightarrow X = 30,\quad Y = 0
\]

所以最优消费组合为 $X^*=30,\ Y^*=0$。


\subsection{完全互补品（里昂惕夫）Leontief}
$U = min\{aX, bY\}, a>0, b>0$ \\
例子： $U = min\{2X,Y\}$:代表1单位的$Y$可以和2单位的$X$互补，形成1个单位的效用。\\
无差异曲线：L型，拐角处效用最大\\
均衡条件：$aX = bY$（即拐角处条件），并且满足预算约束 $P_x \cdot x + P_y \cdot y = I$（不能求导）。\\

\textbf{Example1:}某消费者的效用函数为$U=min\{X,2Y\}$，收入为100元，$P_X=4$，$P_Y=2$。最优消费组合是：\\
完全互补品的最优比例满足$X=2Y$。代入$4X+2Y=100$，得到$Y=10$、$X=20$。


\section{备用}
$\text{对于任意一条无差异曲线：}$
$$U = (Q_1, Q_2) $$ 
$$dU = \frac{\partial U}{\partial Q_1}dQ_1 + \frac{\partial U}{\partial Q_2}dQ_2$$
$$ 0 = MU_1 dQ_1 + MU_2 dQ_2 $$
$$ MRS = -\frac{dQ_2}{dQ_1} = -\frac{MU_1}{MU_2} $$


\end{document}
